\section{Discussion}
\label{sec:discussion}

\subsection{Validating Structural Reasoning}

Five-phase validation provides evidence that the LLM performs structural reasoning rather than pattern matching. The 69.1 percentage point discrimination between 2024 (81.2\%) and 2020 (12.1\%) establishes selectivity: if the model relied on memorized associations or universal statistical patterns, detection rates would converge across periods rather than exhibiting 5.7$\times$ separation.

Three findings support this interpretation. First, the LLM maintained 98\% mechanical accuracy across both detection conditions---correctly computing persistence, magnitude, and flip counts regardless of whether windows qualified as regimes. Second, negative controls achieved 0\% false positives on transitional and low-magnitude synthetic data, confirming the model enforces structural criteria rather than over-detecting patterns. Third, confidence scores provide continuous quality discrimination beyond binary classification ($r > 0.4$ with regime metrics), suggesting the model tracks regime robustness, not just threshold crossing.

This selectivity contrasts with earlier 5-day trajectory testing~\cite{regan2025obfuscation}, which achieved 98--100\% detection across all market conditions by identifying universal daily hedging flows. The deliberate shift to 30-day windows with strict criteria transformed the task from universal pattern matching to selective regime identification---a critical distinction for validating genuine structural reasoning.

\subsection{Market Structure Evolution}

The multi-year temporal analysis (Phase 5) reveals gradual regime evolution tracking 0DTE options adoption rather than sharp structural breaks. Detection progressed from 12.2\% (2020) through borderline years (3.7\% in 2021, 32.4\% in 2022) to sustained 100\% detection in 2024--2025, with GEX magnitude growing 360\% (\$5B $\rightarrow$ \$23B).

\begin{figure}[!htbp]
\centering
\includegraphics[width=\textwidth]{../figures/output/fig08_detection_progression.png}
\caption{Temporal progression of regime detection (2020--2025). Detection rates evolved from pre-regime baseline through growing adoption to structural shift at 100\% in 2024--2025. GEX magnitude grew 360\%, far exceeding inflation.}
\label{fig:progression}
\end{figure}

The non-monotonic pattern (32.4\% in 2022 dipping to 20.2\% in 2023 before reaching 100\% in 2024) is informative: 2023's elevated volatility (FOMC uncertainty, banking stress) repeatedly disrupted regime formation despite growing 0DTE hedging pressure, suggesting regime persistence requires both sustained dealer pressure \textit{and} a volatility environment permitting consolidation. This tipping-point dynamic strengthens the structural interpretation.

While the temporal coincidence with 0DTE proliferation is consistent with a causal relationship, we acknowledge that concurrent factors---interest rate changes (0.25\% $\rightarrow$ 5.5\%), volatility regime shifts, quantitative trading growth, and market maker concentration---may contribute. However, the non-monotonic detection pattern argues against gradual secular trends as primary drivers: detection remained flat through 2023 despite continuous interest rate increases and technology diffusion, then jumped discontinuously in 2024 coinciding with 0DTE market saturation ($\approx$48\% volume share).

\subsection{Temporal Obfuscation as Generalizable Methodology}

Our obfuscation testing framework generalizes beyond financial applications. The core principle---stripping temporal and contextual identifiers to force structural reasoning---addresses a universal challenge in deploying LLMs for domain-specific analysis: distinguishing genuine reasoning from training data contamination. The methodology requires three components: (1) identifiable features that can be systematically removed, (2) structural patterns that persist after obfuscation, and (3) negative controls validating discrimination.

Potential applications include medical time series analysis (removing patient identifiers and dates while preserving physiological patterns), climate data analysis (obfuscating location and temporal markers), and industrial anomaly detection (stripping equipment identifiers). In each case, temporal obfuscation could validate whether LLMs reason from structural properties or rely on memorized associations from training data.

\subsection{Limitations}

Five limitations merit discussion.

\noindent\textbf{Single asset}: Testing focused on SPY; cross-asset generalization (individual equities) remains untested.

\noindent\textbf{End-of-day measurement}: Our GEX\_OI approach captures dealer inventory but not intraday gamma dynamics; high-frequency flow data could refine detection.

\noindent\textbf{Causal attribution}: The 0DTE hypothesis is supported by temporal coincidence and theoretical mechanism but remains circumstantial; observational data cannot exclude alternative explanations.

\noindent\textbf{Framework-dependent reasoning}: Ablation testing (n=84) confirmed identical classification accuracy with and without narrative scaffolding, suggesting the framework's value lies in interpretability rather than detection accuracy.

\noindent\textbf{Threshold sensitivity}: All 25 tested parameter combinations maintained $>$72pp discrimination, but thresholds represent empirically validated design choices rather than first-principles derivations.
